\documentclass[openany,fontset=fandol]{ctexbook}
\usepackage{amsmath,amssymb}
\usepackage[unicode,hidelinks]{hyperref}
\setlength{\emergencystretch}{2em}
\interfootnotelinepenalty=10000
\newcommand{\OO}{\mathcal O}
\newcommand{\LL}{\mathcal L}
\newcommand{\MM}{\mathcal M}
\newcommand{\HH}{\mathcal H om}
\newcommand{\Div}{\operatorname{Div}}
\newcommand{\Pic}{\operatorname{Pic}}
\newcommand{\PP}{\mathbf P}
\begin{document}
\title{代数几何基础第五卷预备笔记\\固定可逆模的相对除子与可表性}
\author{据法文扫描和英译编辑稿互校续译}
\date{2026年11月22日}
\maketitle
\chapter*{范围与接续}
第76批，第三篇第2小节2.11至定理2.13及其证明末；法扫描物理123--125页（手写130--132），英译编辑稿物理95--96页。本批完成第2小节，下一处是第3小节“除子的线性系与到射影丛的态射”。本批尚未并入1158页合订本；正式出版的第四卷第一分册另行接续。
\chapter*{第三篇（续）}
\section*{2　多重射影丛上的可逆层与相对除子（续）}
\noindent\textbf{2.11。}\footnote{法稿此处直接编号2.11，英译编辑稿误标为“Corollary 2.11”。英稿夹有译者关于“抽象化”的按语，不属于原作者的正文。}

前面对 \(\Div_{P/S}\) 的描述之所以如此简单，是因为
\(\Pic(P/S)\) 的结构很简单（确切说，缘于 Picard 预概形
\(\Pic_{P/S}\) 的“离散”结构……）。我们可以把刚才的论证抽象出来，
其要点是建立相对于 \(\Pic\) 的一种可表性。

为此，给定一个固有、平坦且有限表示的态射 \(f:X\to S\)，以及 \(X\) 上的可逆模 \(\LL\)。
我们要确定 \(\Div^+(X/S)\) 的子集 \(\Div^{\LL}(X/S)\)，
它由这样一些相对有效除子 \(D\) 组成：\(\OO_X(D)\) 同构于
\(\LL\otimes_{\OO_X} f^*\MM\)，其中 \(\MM\) 是 \(S\) 上的可逆模，
可以随 \(D\) 而变。\footnote{两稿均称此为“子群”（法文 \emph{sous-groupe}），
并省写了从底的拉回。固定 \(\LL\) 后一般连加法都不封闭：
取 \(X=\PP^1_k\)、\(\LL=\OO(1)\)，一个有理点的次数为1，
两个这样的除子之和次数为2，因而不再属于该子集。原稿在2.13已改用“子函子”表述。}

设单位映射 \(\OO_S\to f_*\OO_X\) 是同构。这样，上面的
\(\MM\) 可由 \(D\) 典范恢复：
\[
 \MM\simeq f_*\bigl(\LL^{-1}\otimes_{\OO_X}\OO_X(D)\bigr).
\]
因此，给定一个与固定 \(\LL\) 相应的 \(D\)，就等价于给出
\(\LL\otimes f^*\MM\) 的一个横截正则截面 \(\phi\)，其中
\(\MM\) 与截面一起按同构识别。\footnote{“横截正则”是指每条纤维上的局部方程都是非零因子，
不是处处不取零值。典范恢复使用投影公式：可逆模局部平凡，
因此 \(f_*f^*\MM\simeq\MM\otimes f_*\OO_X\simeq\MM\)。
唯一性是指与已给同构（或截面）相容的同构，并不是说可逆模没有自同构。}

由第三章第7节，存在 \(S\) 上有限表示模 \(Q\)，其构造与任意基变换交换；
若令 \(G\) 在 \(S\) 上的拟凝聚模中变化，则有函子的同构
\[
 f_*(\LL\otimes_{\OO_X} f^*G)\simeq\HH_{\OO_S}(Q,G).
\]
（准确说，所引之处设 \(S\) 局部诺特；利用这里所用的基变换相容性，
可循通常的诺特逼近方法去掉这一假设。）特别地，给出 \(\phi\)
等价于给出一个同态 \(\psi:Q\to\MM\)。
若 \(\phi\) 横截正则，则它限制到每条纤维上不是零截面；
若 \(X\) 的各纤维都是整概形，则这一条件也是充分的。
以 \(\psi\) 来说，这正是 \(\psi\) 为满射的条件，因此
\(\psi\) 对应于射影丛 \(\PP(Q)\) 在 \(S\) 上的一个截面。\footnote{此处用商约定，
\(\PP(Q)=\operatorname{Proj}_S(\operatorname{Sym}Q)\)；
\(Q\) 只须有限表示，不预设局部自由。基变换相容性在这里用于整个所表示的函子：
到 \(k(s)\) 后的 \(\psi_s\) 非零，当且仅当对应的 \(\phi_s\) 非零。
因 \(\MM_s\) 是一维向量空间，\(\psi_s\ne0\) 等价于 \(\psi_s\) 满射；
再对有限型余核使用 Nakayama 引理。整纤维上的非零截面在泛点非零，
在任一局部平凡化下的芽都是整环中的非零元，所以是非零因子。}
于是得到：

\medskip
\noindent\textbf{命题2.12。}
沿用上述记号。集合 \(\Div^{\LL}(X/S)\) 与 \(\PP(Q)\) 在 \(S\) 上的满足下列条件的截面典范一一对应：
截面所给的可逆商 \(\psi:Q\twoheadrightarrow\MM\) 所定义的
\(\LL\otimes f^*\MM\) 的截面 \(\phi\) 是横截正则的。

现在进一步假设 \(\OO_S\to f_*\OO_X\) 在任意基变换后仍为同构。
由第三章第7节，这等价于对每个 \(s\in S\) 都有自然同构
\begin{equation}
 k(s)\xrightarrow{\;\sim\;} H^0(X_s,\OO_{X_s}).
 \tag{xxx}
\end{equation}
于是，对任意基变换 \(S'\to S\)，命题2.12也适用于 \(X_{S'}/S'\)。
由此得到以下定理。\footnote{普遍成立的是单位映射为同构，不只是原底 \(S\) 上的等式。
现代参照见 Stacks Project \href{https://stacks.math.columbia.edu/tag/0BDM}{第36.32节}，
引理36.32.5与36.32.6。}

\clearpage
\noindent\textbf{定理2.13。}
设 \(f:X\to S\) 固有、平坦、有限表示，并满足上式 \((\mathrm{xxx})\)；
设 \(\LL\) 为 \(X\) 上的可逆模。考虑上文由 \(\LL\) 定义的
\(\Div^+_{X/S}\) 的子函子 \(\Div^{\LL}_{X/S}\)。
存在 \(S\) 上有限表示模 \(Q\)，使此子函子可由射影丛
\(\PP(Q)\) 中一个准紧开浸入的开子概形表示。
若 \(X/S\) 的纤维几何整，则 \(\Div^{\LL}_{X/S}\) 可由整个
\(\PP(Q)\) 表示。\footnote{英译编辑稿定理陈述漏列“平坦”，并有下标错置；
法稿明载 \emph{propre et plat}。几何整性保证纤维经过任意扩域仍为整概形，
这里的“若”不能译成对所有 \(X/S\) 的无条件结论。
法语 \emph{rétrocompact} 指开浸入准紧，不要求基概形本身准紧。}

最后一项由前面关于整纤维的论证直接得到。一般情形中，已有单态射
\[
 \Div^{\LL}_{X/S}\longrightarrow\PP(Q).
\]
只须证明它可由准紧开浸入表示。为此，在任意换底后，
给定 \(\PP(Q)\) 的一个截面即可检验；仍以 \(S\) 记新的底。
这个截面就是 \(Q\) 的一个可逆商 \(\MM\)，从而给出
\(\LL\otimes f^*\MM\) 的截面 \(\phi\)，它在每条纤维上的限制都不是零截面。
在 \((\mathrm{Sch})/S\) 上，终函子 \(S\) 有一个子函子：
它把使 \(\phi\) 变成横截正则的那些换底挑出来。
要证明的就是这个子函子可由 \(S\) 的一个准紧开浸入的开子概形表示。
它是开子概形，因为 \(f\) 固有，横截正则性又是开条件（参见第11节……）；
开浸入的准紧性则归约到诺特情形。\footnote{开性的论证是：
在 \(X\) 中先取横截正则性的开域；其闭补在固有映射 \(f\) 下的像是闭的，
在 \(S\) 中取该闭集的补。纤维上的正则性可沿剩余域扩张检验，
因而得到的开集刻画任意换底，而非只刻画原底上的点。
准紧性另证：在底的每个仿射开集上，把有限表示数据、平坦且固有的
态射、可逆模及截面下降到一个诺特模型；横截正则的条件与基变换相容，
相应开集由诺特底上的开集拉回。可参照 Stacks Project
\href{https://stacks.math.columbia.edu/tag/062Y}{相对有效 Cartier 除子判据}。}
\end{document}

